%%%%%%%%%%%%%%%%%%%%%%%%%%%%%%%%%%%%%%%%%%%%%%%%%%%%%%%%%%%
%% मराठी भाषा आणि साहित्यावर AI चा प्रभाव
%% Marathi Language & Literature — AI Impact Seminar
%% Beamer Presentation (frames only — insert into your master .tex)
%%%%%%%%%%%%%%%%%%%%%%%%%%%%%%%%%%%%%%%%%%%%%%%%%%%%%%%%%%%

%%%%%%%%%%%%%%%%%%%%%%%%%%%%%%%%%%%%%%%%%%%%%%%%%%%%%%%%%%%
%% SLIDE 1 — Opening Hook
%%%%%%%%%%%%%%%%%%%%%%%%%%%%%%%%%%%%%%%%%%%%%%%%%%%%%%%%%%%
\begin{frame}[fragile]\frametitle{ही कविता कोणाची?}

\begin{block}{कविता A}
पाऊस पडला झिम झिम झिम  \\
अंगण झाले ओले चिंब  \\
पाऊस पडला मुसळधार  \\
रान झाले हिरवेगार
\end{block}

\begin{block}{कविता B}
पाऊस आला, आठवणी जाग्या झाल्या  \\
ओल्या मातीमध्ये बालपण सापडलं  \\
छपरावरच्या थेंबांत गाणं ऐकलं  \\
मन नकळत पुन्हा हरवलं
\end{block}

% \vspace{0.3cm}
% {\small एक कविता AI ने तयार केली आहे. एक माणसाने लिहिलेली आहे.}\\
% \textbf{तुम्ही ओळखू शकता का?}

\end{frame}

%%%%%%%%%%%%%%%%%%%%%%%%%%%%%%%%%%%%%%%%%%%%%%%%%%%%%%%%%%%
%% SLIDE 2 — Why This Topic Matters
%%%%%%%%%%%%%%%%%%%%%%%%%%%%%%%%%%%%%%%%%%%%%%%%%%%%%%%%%%%
\begin{frame}[fragile]\frametitle{हा विषय आज का महत्त्वाचा?}

\begin{itemize}
\item कृत्रिम बुद्धिमत्ता - आर्टिफिशिअल इंटेलिजन्स - ए.आय. - AI
\item AI वेगाने दैनंदिन जीवनाचा भाग बनत आहे
\item लेखन, साहित्य, शिक्षण, सर्व क्षेत्रात बदल घडतोय
\item मराठी भाषेसाठी मोठी संधी \textit{आणि} तितकाच धोका
\item तरुण पिढी मोठ्या प्रमाणात AI वापरत आहे
\item साहित्यिक, शिक्षक, धोरणकर्त्यांनी आताच सजग राहणे गरजेचे
\end{itemize}

\end{frame}

%%%%%%%%%%%%%%%%%%%%%%%%%%%%%%%%%%%%%%%%%%%%%%%%%%%%%%%%%%%
%% SLIDE 3, What is AI (Simple)
%%%%%%%%%%%%%%%%%%%%%%%%%%%%%%%%%%%%%%%%%%%%%%%%%%%%%%%%%%%
\begin{frame}[fragile]\frametitle{AI म्हणजे काय? (सोप्या भाषेत)}

\begin{itemize}
\item AI = उदाहरणांमधून शिकणारी यंत्रणा
\item मोबाईलचे predictive typing, हे AI चे छोटे उदाहरण
\item ChatGPT सारखी साधने = मोठे predictive typing
\item माणसासारखा मजकूर तयार करते, पण \textbf{विचार करत नाही}
\item फक्त pattern ओळखतो आणि पुढचा शब्द सुचवतो
\end{itemize}

\end{frame}

%%%%%%%%%%%%%%%%%%%%%%%%%%%%%%%%%%%%%%%%%%%%%%%%%%%%%%%%%%%
%% SLIDE 4, How AI Works
%%%%%%%%%%%%%%%%%%%%%%%%%%%%%%%%%%%%%%%%%%%%%%%%%%%%%%%%%%%
\begin{frame}[fragile]\frametitle{AI कसे काम करते?}

\begin{itemize}
\item लाखो पुस्तके, लेख, वेबसाईट्सवर प्रशिक्षण
\item शब्द व वाक्यांचे संबंध शिकते
\item प्रश्न विचारला की संभाव्य उत्तर तयार करते
\item हे उत्तर ``ज्ञान'' नव्हे, \textit{अंदाज} असतो
\item म्हणून चुकीची माहितीही येऊ शकते
\item \textbf{AI वापरताना स्वतःचा विचार अत्यावश्यक}
\end{itemize}

\end{frame}

%%%%%%%%%%%%%%%%%%%%%%%%%%%%%%%%%%%%%%%%%%%%%%%%%%%%%%%%%%%
%% SLIDE 5, What is a Prompt
%%%%%%%%%%%%%%%%%%%%%%%%%%%%%%%%%%%%%%%%%%%%%%%%%%%%%%%%%%%
\begin{frame}[fragile]\frametitle{Prompt म्हणजे काय?}

\begin{itemize}
\item AI ला दिलेली सूचना म्हणजे \textbf{Prompt}
\item Prompt जितका स्पष्ट तितके उत्तर चांगले
\end{itemize}

\vspace{0.2cm}
\begin{block}{उदाहरण, Prompt कसा सुधारावा}
\begin{enumerate}
\item ``पावसावर कविता लिहा''
\item ``कोकणातील पावसावर भावनिक कविता लिहा''
\item ``कोकणातील पावसावर एका शेतकऱ्याच्या दृष्टिकोनातून कविता लिहा''
\end{enumerate}
\end{block}

\vspace{0.2cm}
\textbf{AI = शिष्य \quad|\quad Prompt = गुरु}

\end{frame}

%%%%%%%%%%%%%%%%%%%%%%%%%%%%%%%%%%%%%%%%%%%%%%%%%%%%%%%%%%%
%% SLIDE 6, Practical Uses of AI
%%%%%%%%%%%%%%%%%%%%%%%%%%%%%%%%%%%%%%%%%%%%%%%%%%%%%%%%%%%
\begin{frame}[fragile]\frametitle{AI चा प्रत्यक्ष वापर, उदाहरणे}

\begin{itemize}
\item विद्यार्थी निबंध आणि प्रकल्प लिहितात
\item शिक्षक प्रश्नपत्रिका व नोट्स तयार करतात
\item पत्रकार बातम्यांचे ड्राफ्ट तयार करतात
\item YouTube व्हिडिओसाठी स्क्रिप्ट लेखन
\item सोशल मीडिया पोस्ट काही सेकंदात तयार
\item लेखक नवीन कल्पना मिळवण्यासाठी AI वापरतात
\end{itemize}

\end{frame}

%%%%%%%%%%%%%%%%%%%%%%%%%%%%%%%%%%%%%%%%%%%%%%%%%%%%%%%%%%%
%% SLIDE 7, AI in Writing
%%%%%%%%%%%%%%%%%%%%%%%%%%%%%%%%%%%%%%%%%%%%%%%%%%%%%%%%%%%
\begin{frame}[fragile]\frametitle{लेखनात AI चे उपयोग}

\begin{itemize}
\item व्याकरण तपासणी आणि Proofreading
\item भाषांतर (मराठी $\leftrightarrow$ इंग्रजी व इतर भाषा)
\item लेख, भाषण, कथा यांचा मसुदा तयार करणे
\item मोठ्या मजकुराचा सारांश
\item विविध शैलीत लेखन, भावनिक, औपचारिक, विनोदी
\end{itemize}

\end{frame}

%%%%%%%%%%%%%%%%%%%%%%%%%%%%%%%%%%%%%%%%%%%%%%%%%%%%%%%%%%%
%% SLIDE 8, AI in Education
%%%%%%%%%%%%%%%%%%%%%%%%%%%%%%%%%%%%%%%%%%%%%%%%%%%%%%%%%%%
\begin{frame}[fragile]\frametitle{शिक्षण क्षेत्रात AI}

\begin{itemize}
\item विद्यार्थ्यांना वैयक्तिक मार्गदर्शन
\item नोट्स व सारांश त्वरित उपलब्ध
\item शिक्षकांचा वेळ वाचवणारे साधन
\item कठीण संकल्पना सोप्या भाषेत समजावणे
\item \textbf{धोका:} विद्यार्थ्यांची स्वतंत्र विचारशक्ती कमी होण्याची शक्यता
\end{itemize}

\end{frame}

%%%%%%%%%%%%%%%%%%%%%%%%%%%%%%%%%%%%%%%%%%%%%%%%%%%%%%%%%%%
%% SLIDE 9, AI for Indian / Marathi Languages
%%%%%%%%%%%%%%%%%%%%%%%%%%%%%%%%%%%%%%%%%%%%%%%%%%%%%%%%%%%
\begin{frame}[fragile]\frametitle{भारतीय भाषांसाठी AI, मराठीवर लक्ष}

\begin{itemize}
\item \textbf{Bhashini}, भारतीय भाषांमधील भाषांतर साधन (केंद्र सरकार)
\item Modi Script OCR, जुन्या मोडी दस्तऐवजांचे डिजिटायझेशन
\item मराठी तांत्रिक शब्दसंग्रह निर्मिती प्रकल्प
\item \textbf{मराठी डेटा कमी}, AI मॉडेल्स प्रशिक्षणासाठी अधिक मजकूर हवा
\item सरकारी व खाजगी उपक्रमांची गरज
\end{itemize}

\end{frame}

%%%%%%%%%%%%%%%%%%%%%%%%%%%%%%%%%%%%%%%%%%%%%%%%%%%%%%%%%%%
%% SLIDE 10, Opportunities for Marathi
%%%%%%%%%%%%%%%%%%%%%%%%%%%%%%%%%%%%%%%%%%%%%%%%%%%%%%%%%%%
\begin{frame}[fragile]\frametitle{मराठीसाठी नवीन संधी}

\begin{itemize}
\item Marathi datasets तयार करणे, AI ला मराठी शिकवणे
\item Apps, websites, सरकारी सेवांचे मराठी localization
\item AI models training मध्ये भाषातज्ज्ञांसाठी रोजगार
\item मराठी ऑडिओबुक, पॉडकास्ट, व्हिडिओ content वाढवण्याची संधी
\item डिजिटल मराठी साहित्य जागतिक वाचकांपर्यंत पोहोचवता येणे
\end{itemize}

\end{frame}

%%%%%%%%%%%%%%%%%%%%%%%%%%%%%%%%%%%%%%%%%%%%%%%%%%%%%%%%%%%
%% SLIDE 11, Impact on Jobs
%%%%%%%%%%%%%%%%%%%%%%%%%%%%%%%%%%%%%%%%%%%%%%%%%%%%%%%%%%%
\begin{frame}[fragile]\frametitle{नोकऱ्यांवर परिणाम}

\begin{itemize}
\item काही पारंपरिक भूमिकांमध्ये बदल अटळ
\item प्रभावित क्षेत्रे: पत्रकारिता, विपणन, शिक्षण, कायदा, content निर्मिती
\item नवीन संधी: AI training, content review, prompt engineering
\item \textbf{Skills बदलणे आवश्यक}, AI सोबत काम करणे शिकणे महत्त्वाचे
\end{itemize}

\end{frame}

%%%%%%%%%%%%%%%%%%%%%%%%%%%%%%%%%%%%%%%%%%%%%%%%%%%%%%%%%%%
%% SLIDE 12, Benefits
%%%%%%%%%%%%%%%%%%%%%%%%%%%%%%%%%%%%%%%%%%%%%%%%%%%%%%%%%%%
\begin{frame}[fragile]\frametitle{AI चे फायदे}

\begin{itemize}
\item लेखन जलद आणि सोपे होते
\item नवीन लेखकांना सुरुवात करणे सोपे
\item भाषांतरामुळे मराठी साहित्याला जागतिक पोहोच
\item ज्ञान सहज उपलब्ध, कोणाही माणसाला, कुठेही
\item दिव्यांग व्यक्तींसाठी भाषा आणि साहित्याचा सहज प्रवेश
\end{itemize}

\end{frame}

%%%%%%%%%%%%%%%%%%%%%%%%%%%%%%%%%%%%%%%%%%%%%%%%%%%%%%%%%%%
%% SLIDE 13, Risks and Challenges
%%%%%%%%%%%%%%%%%%%%%%%%%%%%%%%%%%%%%%%%%%%%%%%%%%%%%%%%%%%
\begin{frame}[fragile]\frametitle{धोके आणि आव्हाने}

\begin{itemize}
\item मौलिकता कमी होण्याचा धोका, सर्व लेखन एकसारखे वाटू लागणे
\item सांस्कृतिक संदर्भ चुकीचे येणे (उदा. मराठी सण, परंपरा)
\item विद्यार्थ्यांची स्वतंत्र विचारशक्ती आणि भाषिक क्षमता क्षीण होणे
\item AI कडून चुकीची किंवा दिशाभूल करणारी माहिती
\item मराठी अस्मिता आणि बोलीभाषांचा ऱ्हास होण्याची शक्यता
\end{itemize}

\end{frame}

%%%%%%%%%%%%%%%%%%%%%%%%%%%%%%%%%%%%%%%%%%%%%%%%%%%%%%%%%%%
%% SLIDE 14, Copyright & Ethics
%%%%%%%%%%%%%%%%%%%%%%%%%%%%%%%%%%%%%%%%%%%%%%%%%%%%%%%%%%%
\begin{frame}[fragile]\frametitle{कॉपीराइट आणि नैतिकता}

\begin{itemize}
\item AI ने तयार केलेल्या मजकुराचा मालक कोण?
\item विद्यमान साहित्यावर AI प्रशिक्षण, लेखकांच्या संमतीशिवाय?
\item साहित्यचौर्य आणि बनावट श्रेयाचा धोका
\item AI-generated content उघड करण्याची नैतिक जबाबदारी
\item \textbf{स्पष्ट कायदे आणि नियम आवश्यक}
\end{itemize}

\end{frame}

%%%%%%%%%%%%%%%%%%%%%%%%%%%%%%%%%%%%%%%%%%%%%%%%%%%%%%%%%%%
%% SLIDE 15, Thought-Provoking Questions
%%%%%%%%%%%%%%%%%%%%%%%%%%%%%%%%%%%%%%%%%%%%%%%%%%%%%%%%%%%
\begin{frame}[fragile]\frametitle{विचार करायला लावणारे प्रश्न}

\begin{itemize}
\item भावना नसताना लिहिलेली कविता खरी कविता असते का?
\item AI ने लिहिलेल्या साहित्याला पुरस्कार द्यावेत का?
\item भविष्यात लेखकाची भूमिका नक्की काय असेल?
\item AI मराठी भाषा समृद्ध करेल की गरीब करेल?
\item तंत्रज्ञान आणि संस्कृती यांचा समतोल कसा साधावा?
\end{itemize}

\end{frame}

%%%%%%%%%%%%%%%%%%%%%%%%%%%%%%%%%%%%%%%%%%%%%%%%%%%%%%%%%%%
%% SLIDE 16, Policy Recommendations
%%%%%%%%%%%%%%%%%%%%%%%%%%%%%%%%%%%%%%%%%%%%%%%%%%%%%%%%%%%
\begin{frame}[fragile]\frametitle{धोरणात्मक दृष्टीकोन}

\begin{itemize}
\item मराठी AI साधनांना सरकारी प्रोत्साहन व निधी
\item जुन्या ग्रंथालयांचे, हस्तलिखितांचे डिजिटायझेशन
\item शालेय व महाविद्यालयीन अभ्यासक्रमात AI साक्षरता
\item लेखकांचे बौद्धिक संपदा हक्क संरक्षण करणारे कायदे
\item जबाबदार व नैतिक AI वापरासाठी मार्गदर्शक तत्त्वे
\end{itemize}

\end{frame}

%%%%%%%%%%%%%%%%%%%%%%%%%%%%%%%%%%%%%%%%%%%%%%%%%%%%%%%%%%%
%% SLIDE 17, Future Direction
%%%%%%%%%%%%%%%%%%%%%%%%%%%%%%%%%%%%%%%%%%%%%%%%%%%%%%%%%%%
\begin{frame}[fragile]\frametitle{भविष्यातील दिशा}

\begin{itemize}
\item \textbf{AI + लेखक = सहकार्य}, स्पर्धा नव्हे
\item AI कल्पना सुचवेल, अनुभव व भावना लेखकाच्याच राहतील
\item मराठी साहित्य भाषांतराद्वारे जागतिक स्तरावर
\item डिजिटल व मल्टीमीडिया मराठी साहित्याचा उदय
\item \textbf{AI हे साधन आहे, साहित्यिक नव्हे}
\end{itemize}

\end{frame}

%%%%%%%%%%%%%%%%%%%%%%%%%%%%%%%%%%%%%%%%%%%%%%%%%%%%%%%%%%%
%% SLIDE 18, Conclusion
%%%%%%%%%%%%%%%%%%%%%%%%%%%%%%%%%%%%%%%%%%%%%%%%%%%%%%%%%%%
\begin{frame}[fragile]\frametitle{समारोप}

\begin{itemize}
\item AI मुळे बदल अटळ आहे, त्यापासून दूर पळणे शक्य नाही
\item साहित्याचा आत्मा मात्र अजूनही मानवी आहे
\item योग्य वापर केल्यास मराठी भाषेसाठी मोठी संधी
\item मराठी भाषा जपणे आणि वाढवणे आपल्या हातात
\end{itemize}

\vspace{0.3cm}
\begin{center}
\textbf{\large तंत्रज्ञान वापरा, पण विचारपूर्वक!}
\end{center}

\vspace{0.4cm}
\begin{center}
{\large \textbf{प्रश्न?}}
\end{center}

\end{frame}
%%%%%%%%%%%%%%%%%%%%%%%%%%%%%%%%%%%%%%%%%%%%%%%%%%%%%%%%%%%